% page 266 of the facsimile (image p-265 of the scan)
% block 157.5 x 246.3 mm ; left margin 8.0 mm
\pgc{-12.700mm}{0.000mm}{
\l{\cel{1}258}
\l{}
\l{\sou{kampesh}o.(\sou{bot.)}\cel{2}Arboro dornoza ek Mexikia, Antili, qua donas}
\l{\cel{3}tintivo reda. - L.\cel{1}\sou{haematoxylon campechianum}. - DEFIRS.}
\l{}
\l{\sou{kano}. (\sou{bot.}) Planto di qua la stipo, rekta, havas tuberi de}
\l{\cel{3}qui spricas folii qui formacas gaino ye sua bazo. - II.}
\l{\cel{3}Planto di qua la stipo esas glata, interne kava e plena}
\l{\cel{3}ye medulo, genero ek la familio \textquotedbl{}graminei\textquotedbl{}. -L.\cel{1}\sou{arund}o,}
\l{\cel{3}\sou{phragmites}. - EFIS.}
\l{}
\l{\sou{kanabo}. (\sou{bot.})\cel{2}Planto herbatra, di qua la stipo donas fila-}
\l{\cel{3}menti ek qui onu facas filo, telo, e c. - L.\cel{1}\sou{canabia}\cel{1}sa-}
\l{\cel{3}\sou{tiva}. - FIS.}
\l{}
\l{\sou{kanabino}. (\sou{zool.)}\cel{2}Speco di mikra pasero griza, amansebla,e}
\l{\cel{3}quan onu povas dresar por ke lu siflez arii. - L.\cel{1}\sou{fringilla}}
\l{\cel{3}\sou{cannabina}. - L.}
\l{}
\l{\sou{kanalo}. I. Quaza rivero artificala, destinita ad igar inter-}
\l{\cel{3}komunikar du baseni, du riveri o rivereti, o du parti ek}
\l{\cel{3}un rivero o rivereto. - II. Aquo-voyo naturala od artifi-}
\l{\cel{3}cala en la enireyo di portuo, o paseyo navigebla qua duktas}
\l{\cel{3}a ta enireyo. - DEFIRS.}
\l{}
\l{\sou{kanali}o. Persono nehonesta, desestimenda. - DEFIRS.}
\l{}
\l{\sou{kanapeo}.\cel{2}Sidilo di qua la dors-apogilo esas tale granda ke,}
\l{\cel{3}sur lu, plura personi povas sidar ensemble, od un persono}
\l{\cel{3}sternar su. - DFIRS.}
\l{}
\l{\sou{kanario}. (\sou{zool.)}\cel{2}Serino de Kanarii. - DEFIRS.}
\l{}
\l{\sou{kancelero}.. Gardero di la stato-siglili. - DEFIRS.}
\l{}
\l{\sou{kancero}. (\sou{patol.}) Nomo vulgara di la tumori qui rodas la tisui}
\l{\cel{3}di ula parti di la korpo, di la organi. - EFIS.}
\l{}
\l{\sou{kande.}\cel{2}En la epoko dum qua ; ye la instanto dum qua...-FIS.}
\l{}
\l{\sou{kandelo}. Mecho ek kotono, cirkondita da sebo, da vaxo o da}
\l{\cel{3}irgaltra materio grasa e kombustebla, qua, kande olu}
\l{\cel{3}esas acendita, konsumesas lente, e donas flamo lumiziva.}
\l{\cel{3}- EFIRS.}
\l{}
\l{\sou{kandelabro.}\cel{2}I. Kandeliero kun plura branchi por la bujii,-}
\l{\cel{3}II. Granda suportilo qua portas un o plura lanterni gasala}
\l{\cel{3}od elektrala por lumizar la voyi publika, la peristili,}
\l{\cel{3}salonon e c. - DEFIRS.}
\l{}
\l{\sou{kandio}. Sukro, kristaligita a peci. - DEFIS.}
\l{}
\l{kandida. Qua esas sincera, pro ke lua pura anmo, psiko, kon-}
\l{\cel{3}cienco havas nulo celenda. - EFIS.}
\l{}
\l{\cel{1}kandidato.- La persono qua su propozas pri ofico o plaso qui}
\l{\cel{3}vakas. - EFISDR.}
}
